Typed decision models such as Jev return a label with a numeric
confidence, and their primitives have been benchmarked on single tasks.
We ask what single-task benchmarks leave open: when a workflow chains
heterogeneous typed decisions, which signals identify its failures, and
can they identify the costly errors for which human review is worthwhile?
We chain Jev's categorical, binary and ordinal primitives into a
dependent bug-triage workflow and evaluate it on Mozilla and Linux kernel
bugs. In pre-registered tests on 500 new bugs, a frozen estimator built
on what the workflow decided ranked wrong actions better than final-step
confidence (AUROC 0.76 vs.\ 0.55). The gap replicated across five models
from four sources, and Jev's estimator transferred unchanged to an
independently trained Jev-style model; the signal was mostly the chosen
action's historical error rate. Yet under the pre-registered costs it did
not yield a cost-effective review policy: at a 10\% review budget no
strategy cost less than no review, though an exploratory sensitivity
analysis shows this reverses for three of five models when wrong accepts
cost twice as much. An oracle shows the expensive errors exist, but the signals
tested could not isolate them, and a case-level score failed a
pre-registered feasibility check (power at most 0.20), so the
confirmatory study was stopped before new data. Predicting which
decisions are wrong is not knowing which are worth correcting.
